% Modality Performance Subsection
\subsection{Modality Performance}
\label{sec:modality-performance}

Detection and forecasting approaches use different biosignal modalities. Detection studies prioritize movement sensors, while forecasting studies emphasize autonomic nervous system indicators.

\subsubsection{Modality Usage}

Accelerometer data is the most prevalent modality appearing in nine of 13 studies (69\%). Six detection studies and three forecasting studies use accelerometry (ACC). Electrodermal activity was captured in five studies (38\%). Electrocardiogram (ECG) or heart rate variance (HRV) data appears in five studies. Photoplethysmography (PPG) appears in four studies. Gyroscope data appears in three studies, all in detection applications.

Multi-modal approaches dominate the evidence base. Nine of thirteen studies (69\%) use multiple sensor types. Four studies use single-modality approaches: \textcite{spahrDeepLearningbasedDetection2025} (ACC only), \textcite{reintjesECGBasedDetectionEpileptic2025} (ECG only), \textcite{borujenyDetectionEpilepticSeizure2013} (ACC only), and \textcite{odeDevelopmentEpilepticSeizure2023} (ECG, HRV only).

\subsubsection{Sensor Location}

Wrist-worn devices appear in eight studies (62\%). Arm placement appears in three studies. Chest placement appears in one study. Thigh placement appears in one study. \textcite{borujenyDetectionEpilepticSeizure2013} used multiple sensor locations (right arm, left arm, left thigh).

Wrist placement was preferred because of patient acceptance and the best availability of consumer devices. The Empatica E4 was used in three studies. The Embrace watch and a Fitbit smartwatch both appear in one study respectively.


\subsubsection{Detection Performance by Modality}

The Single-modality accelerometer approaches achieve strong performance. \textcite{spahrDeepLearningbasedDetection2025} reported 96\% sensitivity with 0.0054/h FPR using only ACC data from the Empatica E4. This was the lowest FPR among all detection studies. \textcite{borujenyDetectionEpilepticSeizure2013} reported 85\% sensitivity using 2D ACC, though with only three patients.

ECG-based detection for focal seizures is currently clinically non-viable due to unacceptable false alarm rates. \textcite{reintjesECGBasedDetectionEpileptic2025} compared three anomaly detection methods on single-lead ECG. The sensitivity ranged from 2.44\% to 82.79\%, while FPR ranged from 0.11/h to 65.62/h depending on method and configuration. The lowest FPR of 0.11/h detected only 2.44\% of seizures. 

In addition GTCS/FBTCS detection using ECG and HRV achieved 74\% sensitivity with 0.85/h FPR  \parencite{odeDevelopmentEpilepticSeizure2023}.

\textcite{wangEpilepticSeizureDetection2025} directly compared their modality performance. Attitude angle features (PITCH and ROLL) outperformed raw accelerometer and gyroscope data for seizure detection. The architecture achieved 56.4\% to 95.3\% sensitivity depending on the feature configuration. This finding suggests that feature engineering can enhance the value of standard movement sensors considerably.

\subsubsection{Forecasting Performance by Modality}

Forecasting studies consistently use autonomic signals. \textcite{nasseriAmbulatorySeizureForecasting2021} combined ACC, PPG, EDA, temperature, and HR to achieve AUC 0.75 (SD 0.15). \textcite{meiselMachineLearningWristband2020} used EDA, PPG, temperature, and ACC from the Empatica E4 to achieve 51.2\% sensitivity with IoC 14.1\% in LOSO validation.

\textcite{vielufSeizureMonitoringCombined2025} combined EDA, ACC, and temperature with diary data. This study demonstrated that 24-hour harmonic patterns in EDA and accelerometry contain predictive information above chance.

Multi-modal approaches are common but not universally superior. Single-modality ACC (\textcite{spahrDeepLearningbasedDetection2025}) achieves the best FPR in the included detection literature. Forecasting consistently requires multi-modal autonomic data, but AUC-ROC plateaus around 0.75 regardless of modality combination.
